\Article{22}
Auch wenn wir \textit{a priori} gewiss sind, dass die Reihen für $V,X,Y,Z$ konvergieren, lässt sich im Voraus doch nichts über die Geschwindigkeit der Konvergenz bestimmen. Wären die Orte der magnetischen Kräfte auf einen mäßig großen Raum um den Erdmittelpunkt beschränkt, oder wären die magnetischen Flüssigkeiten in der Erde so verteilt, dass dies jenem Fall gleichkäme, müssten die Reihen sehr schnell konvergieren. Je weiter sich diese Orte dagegen zur Oberfläche hin erstrecken und je unregelmäßiger die Verteilung ist, desto mehr muss man sich auf langsame Konvergenz einstellen.

Bei der praktischen Anwendung ist absolute Genauigkeit unerreichbar. Man verlangt nur einen den Umständen angemessenen Grad der Annäherung. Je langsamer die Konvergenz ist, desto mehr Glieder muss man berücksichtigen, um einen bestimmten Grad der Genauigkeit zu erreichen.

Nun enthält $P'$ drei Glieder und setzt somit die Kenntnis von drei Koeffizienten $g^{1,0},g^{1,1},h^{1,1}$ voraus. $P''$ erfordert fünf Koeffizienten, $P'''$ sieben, $P''''$ neun usw. Betrachten wir also $P',P'',P'''$ usw. als Größen erster, zweiter, dritter Ordnung usw., so ist klar, dass man, wenn die Rechnung bis zu den Größen der Ordnung $n$ einschließlich reichen soll, die Werte von $nn+2n$ \SourceLeaf{150}Koeffizienten bestimmen muss, also zum Beispiel 24, wenn man bis zur vierten Ordnung gehen will.

\expandafter\def\csname HistoricalNoteText18\endcsname{Am wenigsten nachteilig würden sich diese Umstände bei einem solchen Bestimmungsverfahren auswirken, wenn die acht Punkte ganz symmetrisch auf der Erdoberfläche verteilt waren, das heißt, wenn sie mit den Ecken eines der Erdkugel eingeschriebenen Würfels zusammenfielen oder einer solchen Lage doch sehr nahe kämen.}
Jeder gegebene Wert von $X,Y$ oder $Z$ bei gegebenen Werten von $u$ und $\lambda$ liefert uns eine Gleichung zwischen den Koeffizienten. Sind also an einem Ort alle Elemente des Erdmagnetismus bekannt, so liefern sie drei Gleichungen. Dürfte man annehmen, dass nur die Glieder bis zur vierten Ordnung merklich bleiben, so würden theoretisch die vollständigen Beobachtungen an acht verschiedenen Punkten genügen, um alle nötigen Koeffizienten zu bestimmen. Diese Annahme ist jedoch schwerlich zulässig, und die zufälligen Fehler, die allen Beobachtungen anhaften, könnten zusammen mit der Vernachlässigung der Glieder höherer Ordnung die Ergebnisse der Elimination stark verfälschen\HistoricalNote{18}. Um den schädlichen Einfluss dieser Umstände zu mindern, müsste man erheblich mehr beobachtete Größen als Unbekannte zugrunde legen, und zwar von weit auseinanderliegenden Punkten in allen Teilen der Erde, und daraus die unbekannten Größen nach der Methode der kleinsten Quadrate ableiten. So eintönig die Ausführung einer solchen Arbeit auch wäre, da alle Gleichungen linear sind, so könnte doch der außerordentliche Aufwand, der sich aus der großen Zahl der unbekannten Größen und Gleichungen ergibt, selbst den kühnsten Rechner davon abhalten, die Arbeit schon jetzt in dieser Form in Angriff zu nehmen, zumal eingeschlichene unzuverlässige Beobachtungen oder Rechenfehler das Ergebnis ganz verderben könnten.

